\chapter{Validazione del Clustering: Misure, Teoria e Metodologie}
\label{chap:dm-cluster-validity}

In questo capitolo viene affrontata in modo rigoroso e approfondito la problematica della \textbf{Cluster Validity} (validazione dei cluster), ovvero l'insieme dei fondamenti teorici, delle metriche quantitative e delle metodologie statistiche inferenziali atte a valutare l'affidabilità, la significatività e la bontà delle strutture prodotte dagli algoritmi di raggruppamento~\cite{tan2018, jain1988}.

La trattazione copre analiticamente:
\begin{enumerate}[noitemsep]
    \item I \textbf{fondamenti e le sfide della validazione}: l'ambiguità intrinseca del paradigma non supervisionato, il confronto con la validazione supervisionata, la tendenza degli algoritmi a scoprire cluster spuri nel rumore puramente casuale (dimostrazione su distribuzione uniforme) e le motivazioni metodologiche della valutazione;
    \item La \textbf{tassonomia classica delle misure di validazione}: criteri non supervisionati (\textit{interni}), criteri supervisionati (\textit{esterni}), criteri \textit{relativi} e test di tendenza al raggruppamento (\textit{clustering tendency});
    \item La \textbf{statistica di Hopkins}: formalizzazione matematica della tendenza al raggruppamento, generazione di punti di campionamento stocastico, distanze dai primi vicini e interpretazione probabilistica della divergenza dall'uniformità spaziale;
    \item La \textbf{validazione interna basata su matrici di prossimità}: correlazione di Pearson tra matrice di similarità e matrice di incidenza ideale, ispezione visiva delle matrici di dissimilarità riordinate e comportamento comparativo su $K$-Means, Complete Link, DBSCAN e forme non convesse;
    \item La \textbf{coesione e la separazione dei cluster}: formulazione basata su grafi contro formulazione basata su prototipi, la scomposizione analitica della dispersione totale $\mathrm{TSS} = \mathrm{WSS} + \mathrm{BSS}$, dimostrazione algebrica dettagliata del teorema di conservazione ed esempio numerico passo-passo in dimensione unidimensionale;
    \item Il \textbf{coefficiente di Silhouette} (Rousseeuw~\cite{rousseeuw1987}): definizione geometrica punto per punto ($a(i)$, $b(i)$, $s(i)$), derivazione dei casi limite, calcolo del \textit{Silhouette Score} globale e diagnostica fine tramite profili e grafici di silhouette;
    \item Le \textbf{metriche interne avanzate da laboratorio}: l'indice di Davies-Bouldin ($DB$) e il rapporto di Cali\'nski-Harabasz ($CH$), interpretazione geometrica e trade-off computazionali;
    \item Le \textbf{misure esterne di validazione}: distribuzioni empiriche congiunte $p_{ij}$, formulazione formale dell'entropia condizionata e della purezza pesata, studio di caso completo sul corpus testuale \textit{LA Documents} e matrici di contingenza;
    \item Il \textbf{quadro statistico inferenziale e i modelli nulli}: formulazione dell'ipotesi nulla $H_0$ di casualità uniforme, simulazione Monte Carlo su distribuzioni empiriche, determinazione dei $p$-value e dei valori critici per $\mathrm{SSE}$ e correlazione matriciale;
    \item Il \textbf{laboratorio applicativo in Python}: implementazione degli indici tramite \texttt{scikit-learn}, analisi empirica del dataset del corso, stima del numero ottimale di cluster $K^*$ mediante il metodo del gomito (\textit{knee/elbow method}) e tabulazione incrociata (\textit{crosstab}) con variabili categoriche esterne.
\end{enumerate}

% ====================================================================
% SEZIONE 10.1: FONDAMENTI E SFIDE DELLA CLUSTER VALIDITY
% ====================================================================
\section{Fondamenti e Sfide della Validazione del Clustering}
\label{sec:dm-validity-foundations}

L'analisi dei cluster (\textit{cluster analysis}) costituisce uno dei compiti cardinali dell'apprendimento non supervisionato (\textit{unsupervised machine learning}). L'obiettivo risiede nella partizione di una collezione di oggetti o vettori di attributi $\mathcal{D} = \{\mathbf{x}_1, \dots, \mathbf{x}_N\} \subset \mathbb{R}^d$ in gruppi disgiunti o parzialmente sovrapposti (detti \textit{cluster}), in modo tale che le istanze associate al medesimo gruppo manifestino un'elevata somiglianza reciproca (\textbf{coesione intra-cluster}), mentre elementi appartenenti a gruppi differenti risultino chiaramente distinti (\textbf{separazione inter-cluster})~\cite{tan2018}.

\subsection{Divergenza tra Apprendimento Supervisionato e Non Supervisionato}
Nell'apprendimento supervisionato (come la classificazione o la regressione), la presenza di un'etichetta target fornita a priori (\textit{ground truth}) rende la valutazione del modello oggettiva e diretta: la bontà predittiva viene misurata confrontando le previsioni con la verità empirica attraverso metriche ben consolidate quali \textit{Accuracy}, \textit{Precision}, \textit{Recall}, \textit{F-measure} o l'area sotto la curva ROC (\textit{AUC}).

Al contrario, nel paradigma non supervisionato non è disponibile alcuna etichetta di classe esterna:
\begin{quote}
\textit{``Clustering is in the eye of the beholder.''}
\end{quote}
Non esiste una nozione univoca o universale di cosa costituisca un ``raggruppamento naturale''. La struttura geometrica identificata dipende inscindibilmente dalla nozione di similarità adottata e dalla classe di modelli algoritmici impiegata:
\begin{itemize}[noitemsep]
    \item Per l'algoritmo \textbf{$K$-Means}, un cluster è concepito come un raggruppamento compatto, iper-sferico e convesso focalizzato attorno a un baricentroide medio;
    \item Per \textbf{DBSCAN}, un cluster è una componente connessa ad alta densità separata da regioni a bassa densità o rumore, con estensione morfologica arbitraria;
    \item Per il \textbf{Single Linkage Hierarchical Clustering}, un cluster corrisponde a una catena contigua di punti mutualmente vicini, suscettibile al fenomeno del concatenamento (\textit{chaining effect}).
\end{itemize}

\subsection{Il Problema Critico: Clustering su Dati Puramente Casuali}
La motivazione più urgente per introdurre una teoria rigorosa della validazione risiede nel fatto che \textbf{tutti gli algoritmi di clustering producono partizioni anche quando applicati a rumore puro}.

Si consideri un esperimento concettuale condotto su un dominio bidimensionale continuo $[0, 1] \times [0, 1]$, in cui vengono campionati $N$ punti in modo del tutto indipendente da una distribuzione uniforme bivariata:
\begin{equation}
\label{eq:dm-uniform-noise-generation}
\mathbf{x}_i \sim \mathcal{U}\left([0, 1]^2\right), \quad \forall i \in \{1, \dots, N\}
\end{equation}
Per definizione stocastica, questo insieme di dati non racchiude alcuna struttura a cluster, né attrattori o centri di densità. Tuttavia:
\begin{enumerate}
    \item Se si esegue l'algoritmo $K$-Means imponendo ad esempio $K=3$, la procedura suddivide il dominio stocastico in tre celle di Voronoi compatte, posizionando i centroidi in modo tale da minimizzare la devianza intra-cluster locale e restituendo una soluzione apparentemente plausibile;
    \item Se si esegue un algoritmo gerarchico agglomerativo basato su \textit{Complete Link}, esso produce un dendrogramma completo che, se tagliato a un livello arbitrario, restituisce sotto-insiemi con diametri ben definiti;
    \item Se si applica DBSCAN impostando parametri $(\varepsilon, \mathrm{MinPts})$ non tarati, esso può aggregare fluttuazioni locali della densità di Poisson come se fossero micro-strutture dense.
\end{enumerate}

\begin{alertblock}{Assioma Fondamentale della Validità del Clustering}
Qualsiasi algoritmo di clustering è un ottimizzatore geometrico progettato per aggregare dati. Esso troverà invariabilmente dei cluster, anche nel vuoto informativo. Di conseguenza, nessun risultato di clustering può essere accettato senza aver prima verificato:
\begin{enumerate}[noitemsep]
    \item Se il dataset originale contenga un'effettiva \textbf{tendenza al raggruppamento} rispetto a un'ipotesi nulla di rumore uniforme;
    \item Se la partizione individuata sia \textbf{robusta, statisticamente significativa} e non derivi da fluttuazioni casuali o da bias algoritmici.
\end{enumerate}
\end{alertblock}

\subsection{Finalità Metodologiche della Cluster Validity}
Nel ciclo di vita dell'estrazione di conoscenza (KDD), la validazione dei cluster persegue quattro obiettivi operativi distinti:
\begin{enumerate}
    \item \textbf{Filtraggio del rumore e verifica di tendenza}: Stabilire preliminarmente se il dataset manifesti una predisposizione non casuale all'aggregazione, scartando insiemi privi di segnale geometrico;
    \item \textbf{Confronto tra algoritmi eterogenei}: Valutare quale paradigma (partizionale, gerarchico, density-based o probabilistico) modelli più fedelmente la geometria intrinseca delle osservazioni;
    \item \textbf{Selezione del modello e stima degli iperparametri}: Determinare la parametrizzazione ottima, identificando il numero intrinseco di cluster $K^*$ in $K$-Means o il raggio di vicinato $\varepsilon$ in DBSCAN;
    \item \textbf{Valutazione locale dei singoli cluster}: Esaminare la qualità intrinseca di ciascun gruppo all'interno di una medesima partizione, distinguendo nuclei ad alta coesione e separazione da cluster marginali, spuri o contaminati da rumore.
\end{enumerate}

% ====================================================================
% SEZIONE 10.2: TASSONOMIA DELLE MISURE DI VALIDAZIONE
% ====================================================================
\section{Tassonomia delle Misure di Validazione}
\label{sec:dm-validity-taxonomy}

Seguendo l'impostazione formale di Jain e Dubes~\cite{jain1988} e Tan et al.~\cite{tan2018}, le metriche e le metodologie di validazione si ripartiscono in quattro macro-categorie ben distinte in base alla disponibilità di informazioni ausiliarie e alla finalità analitica.

\begin{figure}[htbp]
\centering
\resizebox{\linewidth}{!}{%
\begin{tikzpicture}[
    boxStyle/.style={
        rectangle, rounded corners=6pt, draw=navy!85!black, fill=navy!6, thick,
        font=\footnotesize, text width=6.2cm, inner sep=8pt, align=left
    },
    headerStyle/.style={
        rectangle, rounded corners=6pt, draw=purple!85!black, fill=purple!12, very thick,
        font=\small, align=center, text width=13.6cm, inner sep=8pt
    },
    arrowStyle/.style={-Stealth, thick, draw=navy!85!black}
]
    % Header
    \node[headerStyle] (top) {
        \textbf{\large TASSONOMIA DELLE MISURE DI VALIDAZIONE DEL CLUSTERING}\\[3pt]
        \footnotesize Suddivisione metodologica fondata sulla disponibilità di informazione ausiliaria e sul criterio di confronto
    };

    % Upper row
    \node[boxStyle, below left=1.2cm and 0.3cm of top.south, anchor=north east] (b1) {
        \textbf{\textcolor{navy}{\normalsize 1. Misure Non Supervisionate (Interne)}}\\[3pt]
        \scriptsize \textbf{Nessuna informazione esterna utilizzata}\\[5pt]
        $\bullet$ Valutano la partizione unicamente dai dati o dalle distanze.\\[2pt]
        $\bullet$ \textit{Coesione intra-cluster}: compattezza e densità interna.\\[2pt]
        $\bullet$ \textit{Separazione inter-cluster}: isolamento reciproco tra gruppi.\\[2pt]
        $\bullet$ \textbf{Metriche}: $\mathrm{SSE}$ / $\mathrm{WSS}$, Silhouette, Davies-Bouldin ($DB$).
    };

    \node[boxStyle, below right=1.2cm and 0.3cm of top.south, anchor=north west] (b2) {
        \textbf{\textcolor{darkgreen}{\normalsize 2. Misure Supervisionate (Esterne)}}\\[3pt]
        \scriptsize \textbf{Confronto con partizione esterna di riferimento}\\[5pt]
        $\bullet$ Misurano la concordanza con etichette reali (\textit{ground truth}).\\[2pt]
        $\bullet$ Non guidano l'ottimizzazione, ma validano l'utilità applicativa.\\[2pt]
        $\bullet$ Utili per verificare l'allineamento a categorie note a priori.\\[2pt]
        $\bullet$ \textbf{Metriche}: Entropia condizionata, Purezza, $\mathrm{ARI}$, $\mathrm{NMI}$.
    };

    % Lower row
    \node[boxStyle, below=1.4cm of b1] (b3) {
        \textbf{\textcolor{orange!85!black}{\normalsize 3. Misure Relative (Model Selection)}}\\[3pt]
        \scriptsize \textbf{Confronto comparativo tra partizioni distinte}\\[5pt]
        $\bullet$ Confrontano soluzioni al variare di iperparametri o algoritmi.\\[2pt]
        $\bullet$ Determinazione del numero intrinseco di cluster ($K^*$).\\[2pt]
        $\bullet$ Identificazione del gomito (\textit{knee/elbow}) o di massimi d'indice.\\[2pt]
        $\bullet$ \textbf{Approcci}: Curva $\mathrm{SSE}(K)$, curva Silhouette $(K)$, criterio $\mathrm{BIC}$.
    };

    \node[boxStyle, below=1.4cm of b2] (b4) {
        \textbf{\textcolor{crimson}{\normalsize 4. Tendenza al Raggruppamento (Tendency)}}\\[3pt]
        \scriptsize \textbf{Test preliminare di non-casualità spaziale}\\[5pt]
        $\bullet$ Verificano se il dataset possieda pattern prima del clustering.\\[2pt]
        $\bullet$ Rigetto dell'ipotesi nulla di rumore casuale uniforme.\\[2pt]
        $\bullet$ Evitano di estrarre cluster spuri in assenza di struttura fisica.\\[2pt]
        $\bullet$ \textbf{Test di riferimento}: Statistica di Hopkins ($H$).
    };

    % Arrows
    \draw[arrowStyle] (top.south) -- ++(0,-0.4) -| (b1.north);
    \draw[arrowStyle] (top.south) -- ++(0,-0.4) -| (b2.north);
    \draw[arrowStyle] (b1.south) -- node[fill=white, inner sep=2pt, font=\scriptsize, text=gray!80!black] {\textit{confronto su K}} (b3.north);
    \draw[arrowStyle] (b2.south) -- node[fill=white, inner sep=2pt, font=\scriptsize, text=gray!80!black] {\textit{verifica a priori}} (b4.north);
\end{tikzpicture}%
}
\caption{Tassonomia delle misure di validazione del clustering suddivise per natura dei dati e finalità d'indagine.}
\label{fig:dm-validity-taxonomy}
\end{figure}

Come riassunto formalmente nella Figura~\ref{fig:dm-validity-taxonomy}:
\begin{enumerate}
    \item \textbf{Misure Interne}: operano nello spazio non supervisionato. Verificano che i cluster siano intrinsecamente densi e ben separati. La metrica tipica include l'errore quadratico intra-cluster ($\mathrm{SSE}$), il coefficiente di silhouette e la coerenza tra matrice di dissimilarità e partizione;
    \item \textbf{Misure Esterne}: operano confrontando l'output dell'algoritmo con etichette esterne fornite da esperti di dominio (ad esempio classi diagnostiche mediche o categorie tematiche di documenti);
    \item \textbf{Misure Relative}: costituiscono un protocollo di confronto seriale per determinare la corretta iper-parametrizzazione (come il numero ideale di cluster $K^*$);
    \item \textbf{Tendenza al Raggruppamento}: costituisce il test preliminare di consistenza geometrica applicato prima di intraprendere qualsiasi procedura di partizione.
\end{enumerate}

% ====================================================================
% SEZIONE 10.3: TENDENZA AL RAGGRUPPAMENTO E STATISTICA DI HOPKINS
% ====================================================================
\section{Tendenza al Raggruppamento e Statistica di Hopkins}
\label{sec:dm-hopkins-statistic}

Prima di spendere risorse computazionali nell'ottimizzazione di algoritmi di clustering, è doveroso verificare se i dati possiedano una struttura non uniforme. Questa problematica prende il nome di \textbf{Clustering Tendency}~\cite{tan2018}. La tecnica quantitativa più celebre e rigorosa è la \textbf{Statistica di Hopkins}~\cite{hopkins1954, jain1988}.

\subsection{Formalizzazione Matematica della Statistica di Hopkins}
Sia dato un insieme di dati multivariato $\mathcal{D} = \{\mathbf{x}_1, \dots, \mathbf{x}_N\} \subset \mathbb{R}^d$. La procedura di Hopkins valuta la discrepanza tra la distribuzione empirica delle distanze tra punti reali e la distribuzione attesa sotto un processo stocastico uniforme.

L'algoritmo si articola nelle seguenti fasi:
\begin{enumerate}
    \item Si estrae casualmente e senza reimmissione un sottoinsieme di $p$ punti reali da $\mathcal{D}$, denotato come:
    \begin{equation}
    \label{eq:dm-hopkins-real-sample}
    \mathcal{X} = \{\mathbf{x}_1^*, \mathbf{x}_2^*, \dots, \mathbf{x}_p^*\} \subset \mathcal{D}, \quad \text{con } p \ll N
    \end{equation}
    Tipicamente si sceglie $p \approx 0.05 \cdot N$ o $p \approx 0.10 \cdot N$ per garantire l'indipendenza statistica dei campioni estratti.
    
    \item Per ogni punto estratto $\mathbf{x}_i^* \in \mathcal{X}$, si calcola la distanza euclidea minima rispetto al suo punto più vicino appartenente al rimanente insieme dei dati $\mathcal{D} \setminus \{\mathbf{x}_i^*\}$:
    \begin{equation}
    \label{eq:dm-hopkins-w-distance}
    w_i = \min_{\mathbf{x} \in \mathcal{D} \setminus \{\mathbf{x}_i^*\}} \|\mathbf{x}_i^* - \mathbf{x}\|_2
    \end{equation}
    
    \item Si costruisce una regione iper-rettangolare di campionamento $\Omega \subset \mathbb{R}^d$ che racchiuda l'iper-volume occupato dal dataset:
    \begin{equation}
    \label{eq:dm-hopkins-bounding-box}
    \Omega = \prod_{j=1}^d \left[\min_{k=1}^N x_{kj}, \ \max_{k=1}^N x_{kj}\right]
    \end{equation}
    All'interno di tale bounding box $\Omega$, si generano $p$ punti sintetici indipendenti distribuiti uniformemente:
    \begin{equation}
    \label{eq:dm-hopkins-synthetic-sample}
    \mathcal{Y} = \{\mathbf{y}_1, \mathbf{y}_2, \dots, \mathbf{y}_p\} \stackrel{\mathrm{i.i.d.}}{\sim} \mathcal{U}(\Omega)
    \end{equation}
    
    \item Per ciascun punto sintetico $\mathbf{y}_i \in \mathcal{Y}$, si calcola la distanza minima rispetto al punto più vicino nell'intero dataset reale $\mathcal{D}$:
    \begin{equation}
    \label{eq:dm-hopkins-u-distance}
    u_i = \min_{\mathbf{x} \in \mathcal{D}} \|\mathbf{y}_i - \mathbf{x}\|_2
    \end{equation}
    
    \item La \textbf{statistica di Hopkins} $H$ è definita come il rapporto normalizzato tra la somma delle distanze dei punti sintetici e la somma aggregata di tutte le distanze calcolate:
    \begin{equation}
    \label{eq:dm-hopkins-formula}
    H = \frac{\sum_{i=1}^p u_i}{\sum_{i=1}^p u_i + \sum_{i=1}^p w_i}
    \end{equation}
    In alternativa, talune formulazioni storiche adottano la somma delle distanze elevate alla potenza $d$-esima (coerente con la misura di volume nello spazio $\mathbb{R}^d$), ma la forma lineare~(\ref{eq:dm-hopkins-formula}) è lo standard ampiamente consolidato nella pratica del data mining.
\end{enumerate}

\subsection{Interpretazione Probabilistica dei Valori di $H$}
Il valore di $H$ è intrinsecamente vincolato nell'intervallo continuo $[0, 1]$. L'interpretazione dei regimi limite è illustrata nella Tabella~\ref{tab:dm-hopkins-interpretation}.

\begin{table}[htbp]
\centering
\small
\renewcommand{\arraystretch}{1.25}
\begin{tabularx}{\textwidth}{c c X}
\toprule
\textbf{Valore di $H$} & \textbf{Comportamento di $u_i$ vs $w_i$} & \textbf{Diagnostica Spaziale e Tendenza al Clustering} \\
\midrule
$H \approx 0.5$ & $\sum u_i \approx \sum w_i$ & \textbf{Dati Uniformemente Distribuiti (Assenza di Cluster)}: le distanze dei punti casuali $\mathbf{y}_i$ dai dati reali sono statisticamente identiche alle distanze tra punti reali adiacenti. Il dataset non possiede cluster naturali. \\
\addlinespace
$H \to 1.0$ & $\sum u_i \gg \sum w_i$ & \textbf{Forte Tendenza al Raggruppamento (Cluster Densi)}: i punti reali risiedono in addensamenti compatti (quindi $w_i \approx 0$). I punti casuali $\mathbf{y}_i$ cadono prevalentemente nelle vaste aree vuote inter-cluster, registrando distanze elevate $u_i$. \\
\addlinespace
$H \to 0.0$ & $\sum u_i \ll \sum w_i$ & \textbf{Dati Regolarmente Spaziati (Griglia Iper-cubica)}: i punti reali si respingono reciprocamente formando una distribuzione periodica o reticolare (massima dispersione uniforme). \\
\bottomrule
\end{tabularx}
\caption{Interpretazione quantitativa della statistica di Hopkins per la verifica della tendenza al clustering.}
\label{tab:dm-hopkins-interpretation}
\end{table}

\begin{remark}[Soglia Decisionale Operativa]
Nella pratica dell'analisi multivariata, si considera che un dataset manifesti una tendenza al clustering statisticamente significativa se la media di ripetute stime indipendenti della statistica di Hopkins soddisfa la condizione:
\begin{equation}
\label{eq:dm-hopkins-threshold}
\bar{H} > 0.70 \quad (\text{preferibilmente } \bar{H} \ge 0.75 \text{ o } 0.80)
\end{equation}
Se $H \in [0.45, 0.55]$, i dati devono essere considerati indistinguibili da rumore casuale uniforme e l'applicazione di algoritmi di partizionamento deve essere sospesa o considerata esplorativa.
\end{remark}

% ====================================================================
% SEZIONE 10.4: VALIDAZIONE INTERNA MEDIANTE MATRICI DI PROSSIMITÀ
% ====================================================================
\section{Validazione Interna mediante Matrici di Prossimità}
\label{sec:dm-proximity-matrix-validity}

Un approccio elegante e generale per valutare la coerenza interna di una partizione consiste nel confrontare direttamente la matrice di prossimità osservata con una matrice ideale che rifletta fedelmente l'assegnazione ai cluster individuata dall'algoritmo~\cite{tan2018}.

\subsection{Matrice di Similarità vs Matrice di Incidenza Ideale}
Sia data una collezione di $N$ oggetti $\mathcal{D} = \{\mathbf{x}_1, \dots, \mathbf{x}_N\}$ e sia $P \in \mathbb{R}^{N \times N}$ la corrispondente \textbf{matrice di prossimità} (o matrice di somiglianza/distanza), dove $P_{ij} = \mathrm{dist}(\mathbf{x}_i, \mathbf{x}_j)$ oppure $P_{ij} = \mathrm{sim}(\mathbf{x}_i, \mathbf{x}_j)$.

Data una partizione $\mathcal{C} = \{\mathcal{C}_1, \dots, \mathcal{C}_K\}$ prodotta da un algoritmo di clustering, si definisce la \textbf{matrice di incidenza ideale} (o matrice di adiacenza ideale) $Y \in \{0, 1\}^{N \times N}$ ponendo:
\begin{equation}
\label{eq:dm-ideal-matrix-def}
Y_{ij} = 
\begin{cases}
1, & \text{se } \mathbf{x}_i, \mathbf{x}_j \in \mathcal{C}_k \text{ (medesimo cluster)}, \\
0, & \text{se } \mathbf{x}_i \in \mathcal{C}_k, \mathbf{x}_j \in \mathcal{C}_l \text{ con } k \ne l \text{ (cluster distinti)}.
\end{cases}
\end{equation}

Se $P_{ij}$ rappresenta una misura di \textit{similarità} (normalizzata in $[0, 1]$), in una partizione ideale e perfetta gli oggetti associati al medesimo cluster dovrebbero avere similarità pari a $1$, mentre coppie appartenenti a cluster differenti dovrebbero registrare similarità pari a $0$. Pertanto, una partizione di elevata qualità mostrerà un forte accordo statistico tra $P$ e $Y$.

\subsection{Coefficiente di Correlazione Matriciale di Pearson}
Poiché le matrici $P$ e $Y$ sono simmetriche ed hanno elementi diagonali costanti o irrilevanti ($P_{ii} = 1$ e $Y_{ii} = 1$), il grado di accordo tra la matrice reale e quella ideale viene quantificato calcolando il \textbf{coefficiente di correlazione di Pearson} esteso su tutte le $M = \frac{N(N-1)}{2}$ coppie non banali $(i, j)$ con $i < j$:
\begin{equation}
\label{eq:dm-matrix-correlation-formula}
\rho(P, Y) = \frac{\sum_{i < j} \left(P_{ij} - \bar{P}\right)\left(Y_{ij} - \bar{Y}\right)}{\sqrt{\sum_{i < j} \left(P_{ij} - \bar{P}\right)^2 \sum_{i < j} \left(Y_{ij} - \bar{Y}\right)^2}}
\end{equation}
dove $\bar{P}$ e $\bar{Y}$ indicano le medie empiriche dei valori fuori dalla diagonale principale:
\begin{equation}
\bar{P} = \frac{2}{N(N-1)} \sum_{i < j} P_{ij}, \qquad \bar{Y} = \frac{2}{N(N-1)} \sum_{i < j} Y_{ij}
\end{equation}

\begin{itemize}[noitemsep]
    \item Se $P$ è una matrice di \textbf{similarità}, valori elevati e positivi di $\rho(P, Y) \to +1$ attestano che la configurazione dei cluster rispecchia fedelmente i gradienti di affinità naturale del dataset;
    \item Se $P$ è una matrice di \textbf{distanza/dissimilarità}, la correlazione attesa per cluster ben separati assumerà valori fortemente negativi $\rho(P, Y) \to -1$.
\end{itemize}

\subsection{Ispezione Visiva delle Matrici di Dissimilarità Riordinate}
Uno strumento qualitativo di straordinaria efficacia diagnostica risiede nella visualizzazione grafica a mappa termica (\textit{heatmap}) della matrice delle distanze, previa \textbf{riordinazione degli indici delle istanze}:
\begin{enumerate}
    \item Si raggruppano gli indici di riga e di colonna in base al cluster di appartenenza: prima tutte le istanze del cluster $\mathcal{C}_1$, poi quelle di $\mathcal{C}_2$, e così via fino a $\mathcal{C}_K$;
    \item All'interno di ciascun blocco di cluster, le istanze possono essere ulteriormente ordinate secondo la loro distanza progressiva dal rispettivo centroide o seguendo l'ordinamento delle foglie di un albero gerarchico.
\end{enumerate}

\begin{figure}[htbp]
\centering
\resizebox{\linewidth}{!}{%
\begin{tikzpicture}[scale=0.95, every node/.style={font=\footnotesize}]
    % Cluster 1 block
    \draw[fill=navy!70, draw=black, thick] (0, 3) rectangle (2, 5);
    \node[white, font=\bfseries] at (1, 4) {$\mathcal{C}_1 \times \mathcal{C}_1$};
    \node[above] at (1, 5.1) {\textbf{Cluster 1}};
    
    % Inter-cluster 1-2
    \draw[fill=gray!15, draw=black] (2, 3) rectangle (4, 5);
    \node[gray!80!black] at (3, 4) {$\mathrm{dist} \gg 0$};
    
    % Inter-cluster 1-3
    \draw[fill=gray!10, draw=black] (4, 3) rectangle (6, 5);
    \node[gray!80!black] at (5, 4) {$\mathrm{dist} \gg 0$};
    
    % Inter-cluster 2-1
    \draw[fill=gray!15, draw=black] (0, 1) rectangle (2, 3);
    \node[gray!80!black] at (1, 2) {$\mathrm{dist} \gg 0$};
    
    % Cluster 2 block
    \draw[fill=navy!70, draw=black, thick] (2, 1) rectangle (4, 3);
    \node[white, font=\bfseries] at (3, 2) {$\mathcal{C}_2 \times \mathcal{C}_2$};
    \node[above] at (3, 5.1) {\textbf{Cluster 2}};
    
    % Inter-cluster 2-3
    \draw[fill=gray!20, draw=black] (4, 1) rectangle (6, 3);
    \node[gray!80!black] at (5, 2) {$\mathrm{dist} \gg 0$};
    
    % Inter-cluster 3-1
    \draw[fill=gray!10, draw=black] (0, -1) rectangle (2, 1);
    \node[gray!80!black] at (1, 0) {$\mathrm{dist} \gg 0$};
    
    % Inter-cluster 3-2
    \draw[fill=gray!20, draw=black] (2, -1) rectangle (4, 1);
    \node[gray!80!black] at (3, 0) {$\mathrm{dist} \gg 0$};
    
    % Cluster 3 block
    \draw[fill=navy!70, draw=black, thick] (4, -1) rectangle (6, 1);
    \node[white, font=\bfseries] at (5, 0) {$\mathcal{C}_3 \times \mathcal{C}_3$};
    \node[above] at (5, 5.1) {\textbf{Cluster 3}};

    % Labels
    \node[left] at (-0.1, 4) {\textbf{Cluster 1}};
    \node[left] at (-0.1, 2) {\textbf{Cluster 2}};
    \node[left] at (-0.1, 0) {\textbf{Cluster 3}};

    % Explanatory legend
    \node[anchor=west, text width=6.5cm, align=left] at (6.8, 2.0) {
        \textbf{Firma Grafica a Blocchi Diagonali:}\\[4pt]
        $\bullet$ \textbf{Blocchi diagonali scuri}: distanze intra-cluster prossime allo zero (elevatissima coesione interna).\\[3pt]
        $\bullet$ \textbf{Regioni fuori diagonale chiare}: distanze inter-cluster marcatamente elevate (separazione netta).\\[3pt]
        $\bullet$ In presenza di \textbf{rumore casuale}, la matrice non mostra alcuna struttura a blocchi, apparendo uniforme e granulare su tutta l'estensione.
    };
\end{tikzpicture}%
}
\caption{Rappresentazione schematica della matrice delle distanze riordinata per cluster: la presenza di marcati blocchi quadrati lungo la diagonale principale evidenzia empiricamente la validità della partizione.}
\label{fig:dm-reordered-matrix-schema}
\end{figure}

\subsection{Sensibilità della Correlazione Matriciale alla Geometria dei Cluster}
Sebbene la correlazione matriciale costituisca uno strumento teorico affascinante, essa risente fortemente dei presupposti geometrici di sfericità:
\begin{itemize}
    \item \textbf{Cluster compatti e sferici}: per partizioni prodotte da $K$-Means o \textit{Complete Link}, le distanze intra-cluster sono omogeneamente basse e le distanze inter-cluster costantemente alte. In questo scenario la correlazione matriciale è massimizzata (ad esempio $\rho \approx 0.85$ o $0.90$);
    \item \textbf{Cluster allungati, complessi o non convessi}: se il dataset presenta cluster concentrici ad anello o forme sinuose continue (gestibili brillantemente da DBSCAN o \textit{Single Linkage}), le distanze euclidee dirette tra punti appartenenti all'estremo opposto del medesimo anello sono ampie, pur appartenendo allo stesso gruppo. La matrice di incidenza $Y$ assegnerà $Y_{ij} = 1$, ma $P_{ij}$ registrerà una distanza considerevole. Ne consegue che la correlazione matriciale $\rho(P, Y)$ risulterà ingannevolmente bassa, pur essendo la partizione geometricamente impeccabile.
\end{itemize}

% ====================================================================
% SEZIONE 10.5: COESIONE E SEPARAZIONE DEI CLUSTER
% ====================================================================
\section{Coesione e Separazione: Teoria ed Equivalenza Algebrica}
\label{sec:dm-cohesion-separation-theory}

Le due proprietà cardinali che definiscono la bontà di un partizionamento non supervisionato sono la \textbf{coesione} (\textit{cohesion}, o compattezza interna) e la \textbf{separazione} (\textit{separation}, o isolamento reciproco)~\cite{tan2018}. Nella letteratura scientifica coesistono due prospettive concettuali fondamentali per quantificarle: la formulazione basata su grafi di prossimità e la formulazione basata su prototipi geometrici.

\subsection{Prospettiva Basata su Grafi vs Basata su Prototipi}
\begin{enumerate}
    \item \textbf{Prospettiva basata su grafi}: Le istanze sono nodi di una rete pesata i cui archi hanno peso pari alla somiglianza $s(\mathbf{x}_i, \mathbf{x}_j)$ o alla distanza $d(\mathbf{x}_i, \mathbf{x}_j)$.
    \begin{itemize}[noitemsep]
        \item La \textit{coesione} di un cluster $\mathcal{C}_k$ è la somma dei pesi di tutti gli archi che connettono coppie interne al gruppo:
        \begin{equation}
        \mathrm{Cohesion}(\mathcal{C}_k) = \sum_{\mathbf{x} \in \mathcal{C}_k, \mathbf{y} \in \mathcal{C}_k} s(\mathbf{x}, \mathbf{y})
        \end{equation}
        \item La \textit{separazione} tra due cluster distinti $\mathcal{C}_k$ e $\mathcal{C}_l$ è il peso aggregato degli archi del taglio (\textit{cut edges}):
        \begin{equation}
        \mathrm{Separation}(\mathcal{C}_k, \mathcal{C}_l) = \sum_{\mathbf{x} \in \mathcal{C}_k, \mathbf{y} \in \mathcal{C}_l} s(\mathbf{x}, \mathbf{y})
        \end{equation}
    \end{itemize}
    
    \item \textbf{Prospettiva basata su prototipi}: Ciascun cluster $\mathcal{C}_k$ è condensato nel proprio rappresentante geometrico (centroide $\mathbf{c}_k$). Coesione e separazione vengono misurate valutando la dispersione quadratica dei punti rispetto ai propri centroidi e la dispersione dei centroidi rispetto al baricentro globale del dataset.
\end{enumerate}

\subsection{Formalizzazione Basata su Prototipi: WSS, BSS e TSS}
Sia dato un dataset $\mathcal{D} = \{\mathbf{x}_1, \dots, \mathbf{x}_N\} \subset \mathbb{R}^d$ partizionato in $K$ cluster disgiunti $\mathcal{C}_1, \dots, \mathcal{C}_K$, con cardinalità $n_k = |\mathcal{C}_k|$ tale che $\sum_{k=1}^K n_k = N$.

Si definisce il \textbf{centroide locale} $\mathbf{c}_k$ di ciascun cluster e il \textbf{centroide globale} $\mathbf{m}$ dell'intero dataset:
\begin{equation}
\label{eq:dm-centroids-definitions}
\mathbf{c}_k = \frac{1}{n_k} \sum_{\mathbf{x} \in \mathcal{C}_k} \mathbf{x}, \qquad \mathbf{m} = \frac{1}{N} \sum_{i=1}^N \mathbf{x}_i = \frac{1}{N} \sum_{k=1}^K n_k \mathbf{c}_k
\end{equation}

Si definiscono formalmente le tre quantità di devianza quadratica:
\begin{enumerate}
    \item \textbf{Coesione Intra-Cluster ($\mathrm{WSS}$ o $\mathrm{SSE}$)}: somma degli errori quadratici interni ai singoli cluster (\textit{Within-Cluster Sum of Squares}):
    \begin{equation}
    \label{eq:dm-wss-definition}
    \mathrm{WSS} = \sum_{k=1}^K \sum_{\mathbf{x} \in \mathcal{C}_k} \|\mathbf{x} - \mathbf{c}_k\|_2^2
    \end{equation}
    
    \item \textbf{Separazione Inter-Cluster ($\mathrm{BSS}$)}: dispersione quadratica dei centroidi di cluster rispetto al baricentro globale $\mathbf{m}$, ponderata per la cardinalità di ciascun gruppo (\textit{Between-Cluster Sum of Squares}):
    \begin{equation}
    \label{eq:dm-bss-definition}
    \mathrm{BSS} = \sum_{k=1}^K n_k \|\mathbf{c}_k - \mathbf{m}\|_2^2
    \end{equation}
    
    \item \textbf{Devianza Totale ($\mathrm{TSS}$)}: somma globale degli scarti quadratici di tutti i vettori rispetto al baricentro globale $\mathbf{m}$ (\textit{Total Sum of Squares}):
    \begin{equation}
    \label{eq:dm-tss-definition}
    \mathrm{TSS} = \sum_{i=1}^N \|\mathbf{x}_i - \mathbf{m}\|_2^2
    \end{equation}
\end{enumerate}

\subsection{Dimostrazione Rigorosa del Teorema di Conservazione della Varianza}
Uno dei risultati cardine dell'analisi multivariata e della cluster validity è il teorema di ortogonalità della scomposizione della devianza totale~\cite{tan2018}.

\begin{theorem}[Conservazione della Devianza Globale]
\label{thm:dm-tss-wss-bss}
Per qualsiasi partizione deterministica di un dataset $\mathcal{D}$ in $K$ cluster disgiunti con centroidi definiti come medie aritmetiche, la somma tra la coesione interna ($\mathrm{WSS}$) e la separazione esterna ($\mathrm{BSS}$) equivale esattamente alla devianza totale dei dati ($\mathrm{TSS}$), la quale è costante e indipendente dall'algoritmo di clustering:
\begin{equation}
\label{eq:dm-tss-conservation}
\mathrm{TSS} = \mathrm{WSS} + \mathrm{BSS}
\end{equation}
\end{theorem}

\begin{proof}
Si consideri la definizione della devianza totale $\mathrm{TSS}$ scomponendo la sommatoria su tutti gli $N$ punti nelle $K$ sommatorie disgiunte sui singoli cluster $\mathcal{C}_k$:
\begin{equation}
\mathrm{TSS} = \sum_{i=1}^N \|\mathbf{x}_i - \mathbf{m}\|_2^2 = \sum_{k=1}^K \sum_{\mathbf{x} \in \mathcal{C}_k} \|\mathbf{x} - \mathbf{m}\|_2^2
\end{equation}
Per ogni punto $\mathbf{x} \in \mathcal{C}_k$, si aggiunga e si sottragga il rispettivo centroide locale $\mathbf{c}_k$:
\begin{equation}
\mathbf{x} - \mathbf{m} = (\mathbf{x} - \mathbf{c}_k) + (\mathbf{c}_k - \mathbf{m})
\end{equation}
Calcolando la norma euclidea al quadrato tramite l'espansione del prodotto scalare interiore:
\begin{equation}
\|\mathbf{x} - \mathbf{m}\|_2^2 = \|(\mathbf{x} - \mathbf{c}_k) + (\mathbf{c}_k - \mathbf{m})\|_2^2 = \|\mathbf{x} - \mathbf{c}_k\|_2^2 + \|\mathbf{c}_k - \mathbf{m}\|_2^2 + 2 (\mathbf{x} - \mathbf{c}_k)^\top (\mathbf{c}_k - \mathbf{m})
\end{equation}
Sostituendo questa identità nell'espressione sommatoria su ciascun cluster $\mathcal{C}_k$:
\begin{equation}
\sum_{\mathbf{x} \in \mathcal{C}_k} \|\mathbf{x} - \mathbf{m}\|_2^2 = \sum_{\mathbf{x} \in \mathcal{C}_k} \|\mathbf{x} - \mathbf{c}_k\|_2^2 + \sum_{\mathbf{x} \in \mathcal{C}_k} \|\mathbf{c}_k - \mathbf{m}\|_2^2 + 2 \sum_{\mathbf{x} \in \mathcal{C}_k} (\mathbf{x} - \mathbf{c}_k)^\top (\mathbf{c}_k - \mathbf{m})
\end{equation}
Si analizzino i tre termini singolarmente:
\begin{enumerate}
    \item Il primo termine coincide con l'errore quadratico interno al cluster $k$, la cui somma su tutti i $k$ fornisce per definizione $\mathrm{WSS}$:
    \begin{equation}
    \sum_{k=1}^K \sum_{\mathbf{x} \in \mathcal{C}_k} \|\mathbf{x} - \mathbf{c}_k\|_2^2 = \mathrm{WSS}
    \end{equation}
    
    \item Nel secondo termine, la quantità $\|\mathbf{c}_k - \mathbf{m}\|_2^2$ è costante rispetto alla variabile di sommatoria $\mathbf{x}$. Di conseguenza, sommando per tutti gli $n_k$ punti appartenenti a $\mathcal{C}_k$:
    \begin{equation}
    \sum_{\mathbf{x} \in \mathcal{C}_k} \|\mathbf{c}_k - \mathbf{m}\|_2^2 = n_k \|\mathbf{c}_k - \mathbf{m}\|_2^2
    \end{equation}
    Sommando su tutti i $K$ cluster si ottiene esattamente la definizione di $\mathrm{BSS}$:
    \begin{equation}
    \sum_{k=1}^K n_k \|\mathbf{c}_k - \mathbf{m}\|_2^2 = \mathrm{BSS}
    \end{equation}
    
    \item Si esamini infine il termine di doppio prodotto incrociato. Poiché il vettore $(\mathbf{c}_k - \mathbf{m})$ non dipende da $\mathbf{x}$, esso può essere fattorizzato all'esterno della sommatoria:
    \begin{equation}
    \sum_{\mathbf{x} \in \mathcal{C}_k} (\mathbf{x} - \mathbf{c}_k)^\top (\mathbf{c}_k - \mathbf{m}) = \left[ \sum_{\mathbf{x} \in \mathcal{C}_k} (\mathbf{x} - \mathbf{c}_k) \right]^\top (\mathbf{c}_k - \mathbf{m})
    \end{equation}
    Tuttavia, per definizione di media aritmetica, il baricentro $\mathbf{c}_k$ azzera identicamente la somma degli scarti dei punti da se stesso:
    \begin{equation}
    \sum_{\mathbf{x} \in \mathcal{C}_k} (\mathbf{x} - \mathbf{c}_k) = \sum_{\mathbf{x} \in \mathcal{C}_k} \mathbf{x} - \sum_{\mathbf{x} \in \mathcal{C}_k} \mathbf{c}_k = n_k \mathbf{c}_k - n_k \mathbf{c}_k = \mathbf{0}
    \end{equation}
    Di conseguenza, il termine di prodotto incrociato si annulla rigorosamente per ogni cluster $k$:
    \begin{equation}
    2 \cdot \mathbf{0}^\top (\mathbf{c}_k - \mathbf{m}) = 0
    \end{equation}
\end{enumerate}
Ricongiungendo i risultati, si ottiene:
\begin{equation}
\mathrm{TSS} = \sum_{k=1}^K \left[ \sum_{\mathbf{x} \in \mathcal{C}_k} \|\mathbf{x} - \mathbf{c}_k\|_2^2 + n_k \|\mathbf{c}_k - \mathbf{m}\|_2^2 \right] = \mathrm{WSS} + \mathrm{BSS}
\end{equation}
completando la dimostrazione.
\end{proof}

\begin{figure}[htbp]
\centering
\resizebox{\linewidth}{!}{%
\begin{tikzpicture}[
    scale=0.95,
    barStyle/.style={rectangle, draw=black, thick, minimum height=0.9cm}
]
    % TSS Bar
    \node[barStyle, fill=purple!20, minimum width=11cm] (tss) at (0, 4.0) {
        \textbf{Devianza Totale Globale: $\mathrm{TSS} = \sum \|\mathbf{x}_i - \mathbf{m}\|^2$ (Costante fissa dei dati)}
    };

    % Double arrows connecting
    \draw[Stealth-Stealth, thick] (-5.5, 2.7) -- (5.5, 2.7);
    \node[above, font=\footnotesize\bfseries] at (0, 2.85) {$\mathrm{TSS} = \mathrm{WSS} + \mathrm{BSS} = \text{Costante Invariante}$};

    % Partition 1: Low K (Suboptimal)
    \node[above, font=\footnotesize\bfseries] at (-3.2, 1.8) {Configurazione a basso $K$ (Cluster dispersi)};
    \node[barStyle, fill=crimson!25, minimum width=7.7cm, anchor=west] (w1) at (-5.5, 1.25) {
        \textbf{WSS (SSE) elevato} (Poca coesione)
    };
    \node[barStyle, fill=darkgreen!25, minimum width=3.3cm, anchor=west] (b1) at ($(w1.east)$) {
        \textbf{BSS ridotto}
    };

    % Partition 2: Optimal K
    \node[above, font=\footnotesize\bfseries] at (-3.2, -0.05) {Configurazione ottima $K^*$ (Bilanciamento intrinseco)};
    \node[barStyle, fill=crimson!25, minimum width=2.8cm, anchor=west] (w2) at (-5.5, -0.6) {
        \textbf{WSS basso}
    };
    \node[barStyle, fill=darkgreen!25, minimum width=8.2cm, anchor=west] (b2) at ($(w2.east)$) {
        \textbf{BSS elevato} (Massima separazione tra centroidi)
    };
\end{tikzpicture}%
}
\caption{Rappresentazione del Teorema di Conservazione della Devianza $\mathrm{TSS} = \mathrm{WSS} + \mathrm{BSS}$: all'aumentare della qualità o della complessità del clustering, la minimizzazione della devianza interna $\mathrm{WSS}$ equivale biunivocamente alla massimizzazione della separazione esterna $\mathrm{BSS}$.}
\label{fig:dm-tss-decomposition-bar}
\end{figure}

\begin{alertblock}{Significato Fondamentale per l'Ottimizzazione}
Il Teorema~\ref{thm:dm-tss-wss-bss} sancisce che:
\begin{equation}
\min \mathrm{WSS} \iff \max \mathrm{BSS}
\end{equation}
Minimizzare la somma degli errori quadratici intra-cluster ($\mathrm{SSE}$) in $K$-Means equivale automaticamente e matematicamente a massimizzare la distanza pesata tra i centroidi dei cluster e il baricentro globale. Non è possibile migliorare la coesione interna senza allontanare i gruppi tra loro.

Tuttavia, $\mathrm{BSS}$ presenta una patologia analoga a $\mathrm{WSS}$: poiché all'aumentare di $K$ la coesione interna $\mathrm{WSS}$ decresce monotonicamente fino ad annullarsi per $K=N$, la separazione $\mathrm{BSS}$ cresce monotonicamente verso il tetto massimo $\mathrm{TSS}$. Pertanto, né $\mathrm{WSS}$ né $\mathrm{BSS}$ possono essere impiegati isolatamente per selezionare il valore ottimo di $K$.
\end{alertblock}

\subsection{Esempio Numerico Didattico Unidimensionale}
Per apprezzare il funzionamento analitico della conservazione di varianza, si consideri il dataset unidimensionale classico riportato nelle dispense di riferimento~\cite{tan2018}:
\begin{equation}
\mathcal{D} = \{1, 2, 4, 5\}, \quad N = 4
\end{equation}

\subsubsection{Passo 0: Calcolo del Baricentro Globale e di TSS}
La media globale $\mathbf{m}$ vale:
\begin{equation}
\mathbf{m} = \frac{1 + 2 + 4 + 5}{4} = \frac{12}{4} = 3
\end{equation}
La devianza totale $\mathrm{TSS}$ del dataset vale:
\begin{align}
\mathrm{TSS} &= (1 - 3)^2 + (2 - 3)^2 + (4 - 3)^2 + (5 - 3)^2 \notag \\
&= (-2)^2 + (-1)^2 + 1^2 + 2^2 = 4 + 1 + 1 + 4 = 10
\end{align}
Questo valore costituisce l'invariante geometrico del sistema.

\subsubsection{Caso A: Partizione a Singolo Cluster ($K=1$)}
Se tutti i punti sono assegnati a un unico macro-cluster $\mathcal{C}_1 = \{1, 2, 4, 5\}$:
\begin{itemize}[noitemsep]
    \item Il centroide locale coincide con la media globale: $\mathbf{c}_1 = 3$;
    \item La coesione interna coincide con l'intera devianza:
    \begin{equation}
    \mathrm{WSS} = (1-3)^2 + (2-3)^2 + (4-3)^2 + (5-3)^2 = 10
    \end{equation}
    \item La separazione esterna è nulla, poiché il centroide coincide col baricentro:
    \begin{equation}
    \mathrm{BSS} = 4 \cdot (3 - 3)^2 = 0
    \end{equation}
\end{itemize}
Verifica: $\mathrm{TSS} = 10 + 0 = 10$.

\subsubsection{Caso B: Partizione Naturale a Due Cluster ($K=2$)}
Si partizioni il dataset nei due raggruppamenti naturali intuitivi:
\begin{equation}
\mathcal{C}_1 = \{1, 2\}, \qquad \mathcal{C}_2 = \{4, 5\}
\end{equation}
con cardinalità $n_1 = 2$ e $n_2 = 2$.
\begin{enumerate}
    \item Si calcolano i rispettivi centroidi locali:
    \begin{equation}
    \mathbf{c}_1 = \frac{1 + 2}{2} = 1.5, \qquad \mathbf{c}_2 = \frac{4 + 5}{2} = 4.5
    \end{equation}
    
    \item Si calcola la coesione intra-cluster ($\mathrm{WSS}$):
    \begin{align}
    \mathrm{SSE}(\mathcal{C}_1) &= (1 - 1.5)^2 + (2 - 1.5)^2 = 0.25 + 0.25 = 0.50 \\
    \mathrm{SSE}(\mathcal{C}_2) &= (4 - 4.5)^2 + (5 - 4.5)^2 = 0.25 + 0.25 = 0.50 \\
    \mathrm{WSS} &= \mathrm{SSE}(\mathcal{C}_1) + \mathrm{SSE}(\mathcal{C}_2) = 0.50 + 0.50 = 1.00
    \end{align}
    
    \item Si calcola la separazione inter-cluster ($\mathrm{BSS}$):
    \begin{align}
    \mathrm{BSS} &= n_1 (\mathbf{c}_1 - \mathbf{m})^2 + n_2 (\mathbf{c}_2 - \mathbf{m})^2 \\
    &= 2 \cdot (1.5 - 3)^2 + 2 \cdot (4.5 - 3)^2 \\
    &= 2 \cdot (-1.5)^2 + 2 \cdot (1.5)^2 = 2 \cdot 2.25 + 2 \cdot 2.25 = 4.50 + 4.50 = 9.00
    \end{align}
\end{enumerate}
Verifica del teorema di conservazione:
\begin{equation}
\mathrm{WSS} + \mathrm{BSS} = 1.00 + 9.00 = 10.00 = \mathrm{TSS}
\end{equation}
L'introduzione della partizione corretta a due cluster ha abbattuto la devianza interna del $90\%$ (da $10$ a $1$), convertendola integralmente in separazione tra gruppi ($\mathrm{BSS} = 9$).

% ====================================================================
% SEZIONE 10.6: IL COEFFICIENTE DI SILHOUETTE
% ====================================================================
\section{Il Coefficiente di Silhouette}
\label{sec:dm-silhouette-coefficient}

Introdotto originariamente da Peter Rousseeuw nel 1987~\cite{rousseeuw1987}, il \textbf{Coefficiente di Silhouette} è una delle metriche di validazione interna più versatili ed ampiamente utilizzate nel machine learning. Il suo pregio fondamentale risiede nella capacità di fondere simultaneamente la nozione di coesione interna e di separazione esterna a livello di \textbf{singolo punto dati}, senza dipendere esplicitamente dalla forma sferica o dal baricentroide medio~\cite{tan2018}.

\subsection{Definizione Matematica Punto per Punto}
Sia data una partizione di un dataset $\mathcal{D}$ in $K \ge 2$ cluster $\mathcal{C}_1, \dots, \mathcal{C}_K$. Per ogni specifica istanza $\mathbf{x}_i \in \mathcal{C}_k$:

\begin{enumerate}
    \item Si calcola la \textbf{dissimilarità intra-cluster media} $a(i)$, che esprime la distanza media di $\mathbf{x}_i$ da tutti gli altri punti appartenenti al proprio cluster:
    \begin{equation}
    \label{eq:dm-silhouette-a}
    a(i) = \frac{1}{|\mathcal{C}_k| - 1} \sum_{\mathbf{x}_j \in \mathcal{C}_k, j \ne i} \mathrm{dist}(\mathbf{x}_i, \mathbf{x}_j)
    \end{equation}
    Il parametro $a(i)$ quantifica la \textbf{coesione}: valori piccoli indicano che il punto si trova profondamente integrato e vicino ai compagni di cluster.
    
    \item Si determina la \textbf{distanza media verso i cluster vicini}: per ogni altro cluster $\mathcal{C}_l \ne \mathcal{C}_k$, si calcola la distanza media da $\mathbf{x}_i$ a tutti i punti di $\mathcal{C}_l$:
    \begin{equation}
    d(\mathbf{x}_i, \mathcal{C}_l) = \frac{1}{|\mathcal{C}_l|} \sum_{\mathbf{x}_j \in \mathcal{C}_l} \mathrm{dist}(\mathbf{x}_i, \mathbf{x}_j)
    \end{equation}
    Si definisce $b(i)$ come il minimo di tali medie tra tutti i cluster concorrenti:
    \begin{equation}
    \label{eq:dm-silhouette-b}
    b(i) = \min_{l \ne k} d(\mathbf{x}_i, \mathcal{C}_l)
    \end{equation}
    Il cluster che realizza il minimo in~(\ref{eq:dm-silhouette-b}) viene designato come il \textbf{cluster più vicino} (\textit{neighboring cluster}) per l'istanza $\mathbf{x}_i$. La quantità $b(i)$ quantifica la \textbf{separazione}: valori ampi indicano un marcato isolamento dal cluster contiguo.
    
    \item Il \textbf{coefficiente di Silhouette individuale} $s(i)$ per il punto $\mathbf{x}_i$ è definito come la differenza normalizzata:
    \begin{equation}
    \label{eq:dm-silhouette-s}
    s(i) = \frac{b(i) - a(i)}{\max\left(a(i), \ b(i)\right)}
    \end{equation}
\end{enumerate}

\begin{figure}[htbp]
\centering
\resizebox{\linewidth}{!}{%
\begin{tikzpicture}[
    scale=0.92,
    every node/.style={font=\footnotesize}
]
    % Central formula note at the top
    \node[rectangle, rounded corners=6pt, draw=purple!85!black, fill=purple!8, inner sep=6pt, align=center, font=\small] (formula) at (0, 3.8) {
        $\displaystyle s(i) = \frac{b(i) - a(i)}{\max(a(i), \ b(i))}$\\[3pt]
        \scriptsize $-1 \le s(i) \le 1$
    };

    % Cluster A (own cluster) on the left
    \draw[fill=navy!8, draw=navy!85!black, very thick] (-4.0, 0) ellipse (2.9cm and 1.8cm);
    \node[navy!90!black, font=\bfseries\normalsize] at (-4.0, 2.2) {Cluster di Appartenenza $\mathcal{C}_k$};
    
    % Points in Cluster A
    \filldraw[navy!80!black] (-5.8, 0.4) circle (2.5pt);
    \filldraw[navy!80!black] (-5.0, -1.0) circle (2.5pt);
    \filldraw[navy!80!black] (-5.4, -0.3) circle (2.5pt);
    \filldraw[navy!80!black] (-4.3, 0.8) circle (2.5pt);
    \filldraw[navy!80!black] (-4.0, -0.7) circle (2.5pt);

    % Reference sample point x_i on boundary of cluster A
    \coordinate (xi) at (-2.3, 0.1);
    \filldraw[crimson] (xi) circle (3.5pt) node[above=2pt, crimson, font=\bfseries\normalsize] {$\mathbf{x}_i$};

    % Intra-cluster distance dashed lines
    \draw[dashed, navy!80!black, thick] (xi) -- (-5.8, 0.4);
    \draw[dashed, navy!80!black, thick] (xi) -- (-5.0, -1.0);
    \draw[dashed, navy!80!black, thick] (xi) -- (-4.3, 0.8);
    \draw[dashed, navy!80!black, thick] (xi) -- (-4.0, -0.7);

    % Callout for a(i) positioned below cluster A
    \node[rectangle, rounded corners=5pt, draw=navy!80!black, fill=navy!6, text width=5.6cm, inner sep=6pt, align=center, font=\footnotesize] (boxA) at (-4.0, -2.8) {
        \textbf{Coesione intra-cluster $a(i)$}\\[2pt]
        Media delle distanze da $\mathbf{x}_i$ a tutti gli altri punti del cluster $\mathcal{C}_k$
    };
    \draw[-Stealth, navy!80!black, thick] (boxA.north) -- (-4.0, -1.9);

    % Cluster B (nearest neighboring cluster) on the right
    \draw[fill=darkgreen!8, draw=darkgreen!85!black, very thick] (4.0, 0) ellipse (2.9cm and 1.9cm);
    \node[darkgreen!90!black, font=\bfseries\normalsize] at (4.0, 2.3) {Cluster Concorrente Più Vicino $\mathcal{C}_l$};

    % Points in Cluster B
    \filldraw[darkgreen!80!black] (2.3, 0.5) circle (2.5pt);
    \filldraw[darkgreen!80!black] (2.7, -0.8) circle (2.5pt);
    \filldraw[darkgreen!80!black] (3.8, 0.9) circle (2.5pt);
    \filldraw[darkgreen!80!black] (4.5, -0.5) circle (2.5pt);
    \filldraw[darkgreen!80!black] (5.6, 0.2) circle (2.5pt);

    % Inter-cluster distance dashed lines from x_i
    \draw[dashed, darkgreen!80!black, thick] (xi) -- (2.3, 0.5);
    \draw[dashed, darkgreen!80!black, thick] (xi) -- (2.7, -0.8);
    \draw[dashed, darkgreen!80!black, thick] (xi) -- (4.5, -0.5);

    % Callout for b(i) positioned below cluster B
    \node[rectangle, rounded corners=5pt, draw=darkgreen!80!black, fill=darkgreen!6, text width=5.6cm, inner sep=6pt, align=center, font=\footnotesize] (boxB) at (4.0, -2.8) {
        \textbf{Separazione inter-cluster $b(i)$}\\[2pt]
        Minima distanza media da $\mathbf{x}_i$ verso i punti di un cluster esterno $\mathcal{C}_l$
    };
    \draw[-Stealth, darkgreen!80!black, thick] (boxB.north) -- (4.0, -1.9);
\end{tikzpicture}%
}
\caption{Geometria del coefficiente di Silhouette: confronto visivo tra la coesione interna $a(i)$ e la separazione dal cluster contiguo più vicino $b(i)$.}
\label{fig:dm-silhouette-geometry}
\end{figure}

\subsection{Analisi dei Regimi Limite del Coefficiente}
Dalla presenza del denominatore normalizzante $\max(a(i), b(i))$, scaturisce direttamente che:
\begin{equation}
-1 \le s(i) \le 1, \quad \forall i
\end{equation}
L'interpretazione analitica dei tre regimi estremi è riassunta di seguito:
\begin{itemize}
    \item \textbf{Condizione Ottimale ($s(i) \to +1$)}: Si verifica quando $a(i) \ll b(i)$, ovvero la distanza interna è trascurabile rispetto alla separazione dal cluster contiguo. Il punto risiede nel cuore del proprio raggruppamento, perfettamente isolato dall'esterno.
    \item \textbf{Punto di Confine ($s(i) \approx 0$)}: Si verifica quando $a(i) \approx b(i)$. Il punto giace esattamente sulla frontiera decisionale di Voronoi che separa il cluster $\mathcal{C}_k$ dal cluster contiguo $\mathcal{C}_l$. L'assegnazione ad uno o all'altro gruppo risulta ambigua ed arbitraria.
    \item \textbf{Errore Grave di Assegnazione ($s(i) \to -1$)}: Si verifica quando $a(i) \gg b(i)$. In questo caso, il punto si trova mediamente molto più vicino agli elementi del cluster concorrente rispetto ai membri del gruppo a cui è stato formalmente assegnato dall'algoritmo. Tale circostanza evidenzia un evidente fallimento dell'ottimizzazione o un vincolo rigido sub-ottimale.
\end{itemize}

\subsection{Silhouette Score Globale e Grafici di Profilo}
Per passare dalla valutazione puntuale alla validazione globale dell'intera partizione, si calcola il \textbf{Silhouette Score Medio} $\bar{s}$:
\begin{equation}
\label{eq:dm-mean-silhouette-score}
\bar{s} = \frac{1}{N} \sum_{i=1}^N s(i)
\end{equation}
Un Silhouette Score globale $\bar{s} > 0.70$ denota una struttura eccellente a cluster forti e separati; valori nell'intervallo $[0.50, 0.70]$ attestano una struttura ragionevole; punteggi inferiori a $0.25$ denotano assenza di cluster o forte sovrapposizione.

Oltre al valore medio aggregato, la silhouette consente la costruzione del \textbf{Silhouette Plot} (grafico dei profili):
\begin{enumerate}
    \item Si ordinano le istanze di ciascun cluster in senso decrescente in base al loro coefficiente $s(i)$;
    \item Si traccia un istogramma a barre orizzontali per ogni cluster, visualizzando lo spessore (cardinalità del gruppo) e la penetrazione verso destra lungo l'asse orizzontale $[-1, 1]$;
    \item Si traccia una linea verticale tratteggiata in corrispondenza del valore medio globale $\bar{s}$.
\end{enumerate}
Questo grafico consente di diagnosticare istantaneamente se determinati cluster siano intrinsecamente deboli (ossia contengano numerose barre corte o negative) pur a fronte di un punteggio medio globale apparentemente soddisfacente.

% ====================================================================
% SEZIONE 10.7: INDICI AGGIUNTIVI DA LABORATORIO
% ====================================================================
\section[Metriche Interne: Davies-Bouldin e Cali\'nski-Harabasz]{Metriche Interne: Davies-Bouldin\\e Cali\'nski-Harabasz}
\label{sec:dm-additional-internal-indices}

Nel laboratorio analitico di data mining (in particolare nel notebook sperimentale di clustering~\cite{tan2018}), vengono impiegate metriche complementari che offrono un compromesso computazionale superiore rispetto alla Silhouette (la quale richiede la memorizzazione o il calcolo quadratico delle distanze a coppie $\mathcal{O}(N^2)$).

\subsection{Indice di Davies-Bouldin (DB Score)}
Introdotto da Davies e Bouldin nel 1979~\cite{davies1979}, l'indice $DB$ valuta il peggior rapporto di similarità tra ciascun cluster e il suo rivale più prossimo.

Siano $\mathbf{c}_i$ e $\mathbf{c}_j$ i centroidi dei cluster $\mathcal{C}_i$ e $\mathcal{C}_j$. Si definisce la dispersione intra-cluster media $\sigma_i$ come la distanza media dei punti appartenenti a $\mathcal{C}_i$ rispetto a $\mathbf{c}_i$:
\begin{equation}
\label{eq:dm-db-sigma}
\sigma_i = \frac{1}{|\mathcal{C}_i|} \sum_{\mathbf{x} \in \mathcal{C}_i} \|\mathbf{x} - \mathbf{c}_i\|_2
\end{equation}
La similarità $R_{ij}$ tra la coppia di cluster $\mathcal{C}_i$ e $\mathcal{C}_j$ è definita come:
\begin{equation}
\label{eq:dm-db-rij}
R_{ij} = \frac{\sigma_i + \sigma_j}{d(\mathbf{c}_i, \mathbf{c}_j)}
\end{equation}
dove al numeratore compare la somma delle dispersioni interne (richiesta di coesione) e al denominatore compare la distanza euclidea tra i due centroidi $d(\mathbf{c}_i, \mathbf{c}_j) = \|\mathbf{c}_i - \mathbf{c}_j\|_2$ (richiesta di separazione).

L'indice di \textbf{Davies-Bouldin} $DB$ è definito calcolando la media aritmetica della massima similarità registrata da ciascun cluster:
\begin{equation}
\label{eq:dm-db-score}
DB = \frac{1}{K} \sum_{i=1}^K \max_{j \ne i} R_{ij}
\end{equation}

\begin{property}[Proprietà dell'Indice $DB$]
\begin{itemize}[noitemsep]
    \item Il punteggio $DB$ è limitato inferiormente da zero: $DB \in [0, +\infty)$;
    \item \textbf{Criterio di Ottimalità}: Valori inferiori denotano un clustering di qualità superiore (cluster compatti con dispersioni $\sigma_i$ piccole e ben distanziati con denominatore $d(\mathbf{c}_i, \mathbf{c}_j)$ grande);
    \item Complessità computazionale: calcolabile in tempo $\mathcal{O}(N \cdot d + K^2 \cdot d)$, risultando drasticamente più veloce della Silhouette sui dataset di grandi dimensioni.
\end{itemize}
\end{property}

\subsection{Indice di Cali\'nski-Harabasz (Variance Ratio Criterion)}
Proposto da Cali\'nski e Harabasz nel 1974~\cite{calinski1974}, l'indice $CH$ (detto anche \textit{Variance Ratio Criterion}) rapporta direttamente la separazione inter-cluster $\mathrm{BSS}$ e la coesione intra-cluster $\mathrm{WSS}$, normalizzandole per i rispettivi gradi di libertà:
\begin{equation}
\label{eq:dm-calinski-harabasz}
CH = \frac{\mathrm{BSS} / (K - 1)}{\mathrm{WSS} / (N - K)}
\end{equation}
Valori più elevati di $CH$ indicano cluster maggiormente densi e separati. In corrispondenza del valore intrinseco di $K^*$, la curva $CH(K)$ manifesta frequentemente un picco locale pronunciato.

% ====================================================================
% SEZIONE 10.8: MISURE ESTERNE: ENTROPIA, PUREZZA E CONCORDANZA
% ====================================================================
\section{Misure Esterne: Entropia, Purezza e Corrispondenza}
\label{sec:dm-external-measures}

Quando è disponibile una classificazione esterna indipendente (etichette fornite a priori da esperti di dominio), le metriche supervisionate consentono di misurare fino a che punto i cluster scoperti ricostruiscano o riflettano fedelmente tali categorie~\cite{tan2018}.

\subsection{Distribuzioni di Probabilità Empiriche e Matrice di Contingenza}
Sia data una partizione prodotta dall'algoritmo $\mathcal{C} = \{\mathcal{C}_1, \dots, \mathcal{C}_K\}$ e sia $\mathcal{L} = \{\mathcal{L}_1, \dots, \mathcal{L}_L\}$ la partizione reale delle etichette esterne (\textit{ground truth}), con dimensione totale del dataset pari a $N$.

Si costruisce la \textbf{matrice di contingenza} (o di co-occorrenza) la cui cella generica $m_{ij}$ rappresenta il numero di elementi assegnati contemporaneamente al cluster $\mathcal{C}_i$ e alla classe esterna $\mathcal{L}_j$:
\begin{equation}
\label{eq:dm-contingency-cell}
m_{ij} = |\mathcal{C}_i \cap \mathcal{L}_j|, \qquad m_i = |\mathcal{C}_i| = \sum_{j=1}^L m_{ij}, \qquad m_{\cdot j} = |\mathcal{L}_j| = \sum_{i=1}^K m_{ij}
\end{equation}

Per ciascun cluster $\mathcal{C}_i$, si definisce la distribuzione di probabilità empirica condizionata che un elemento estratto a caso da $\mathcal{C}_i$ appartenga alla classe $\mathcal{L}_j$:
\begin{equation}
\label{eq:dm-empirical-prob}
p_{ij} = \frac{m_{ij}}{m_i} = \frac{|\mathcal{C}_i \cap \mathcal{L}_j|}{|\mathcal{C}_i|}
\end{equation}
ovviamente con $\sum_{j=1}^L p_{ij} = 1$ per ogni cluster fisso $i$.

\subsection{Formulazione Matematica di Entropia e Purezza}
\begin{enumerate}
    \item \textbf{Entropia Individuale di un Cluster ($e_i$)}: Applicando la definizione fondazionale di Entropia di Shannon formalizzata nel Capitolo~\ref{chap:dm-similarity-clustering-foundations} (\autoref{sec:dm-information-theory}), si quantifica il disordine informativo o l'eterogeneità interna del cluster $\mathcal{C}_i$ rispetto alla distribuzione delle classi di riferimento:
    \begin{equation}
    \label{eq:dm-entropy-single}
    e_i = - \sum_{j=1}^L p_{ij} \log_2(p_{ij})
    \end{equation}
    convenendo che $0 \log_2(0) = 0$. L'entropia è minima ($e_i = 0$) se il cluster è perfettamente puro (contiene istanze appartenenti a una singola classe); è massima ($\log_2 L$) se tutte le classi sono equiprobabili nel cluster.
    
    \item \textbf{Entropia Globale Ponderata del Clustering ($E$)}: Media pesata delle entropie individuali, proporzionale alla dimensione dei singoli cluster:
    \begin{equation}
    \label{eq:dm-entropy-global}
    E = \sum_{i=1}^K \frac{m_i}{N} e_i = - \sum_{i=1}^K \frac{m_i}{N} \sum_{j=1}^L p_{ij} \log_2(p_{ij})
    \end{equation}
    Minori valori di $E$ indicano partizioni di qualità superiore rispetto alle etichette di classe.
    
    \item \textbf{Purezza Individuale di un Cluster ($p_i$)}: Rappresenta la frazione del cluster occupata dalla classe maggioritaria in esso contenuta:
    \begin{equation}
    \label{eq:dm-purity-single}
    p_i = \max_j p_{ij} = \max_j \frac{m_{ij}}{m_i}
    \end{equation}
    
    \item \textbf{Purezza Globale Ponderata del Clustering ($P$)}: Media pesata delle purezze individuali:
    \begin{equation}
    \label{eq:dm-purity-global}
    P = \sum_{i=1}^K \frac{m_i}{N} p_i = \frac{1}{N} \sum_{i=1}^K \max_j m_{ij}
    \end{equation}
    La purezza varia nell'intervallo $[0, 1]$: valori prossimi ad $1$ attestano che ciascun cluster è dominato quasi esclusivamente da un'unica classe target.
\end{enumerate}

\subsection{Studio di Caso Analitico: Il Corpus Testuale ``LA Documents''}
Nelle lezioni di riferimento~\cite{tan2018} viene esaminato un esperimento applicativo condotto sul corpus testuale \textit{Los Angeles Times Documents} ($N = 3\,204$ articoli ripartiti su 6 argomenti di cronaca: intrattenimento, finanza, esteri, politica, sport, tecnologia).

Il dataset viene partizionato mediante $K$-Means impostando $K=9$. La matrice di contingenza risultante è riportata integralmente nella Tabella~\ref{tab:dm-la-documents}.

\begin{table}[htbp]
\centering
\small
\renewcommand{\arraystretch}{1.2}
\resizebox{\linewidth}{!}{%
\begin{tabular}{l *{6}{c} c c c}
\toprule
\textbf{Cluster} & \textbf{Ent.} & \textbf{Fin.} & \textbf{Est.} & \textbf{Pol.} & \textbf{Spo.} & \textbf{Tec.} & \textbf{Totale} & \textbf{Purezza} & \textbf{Entropia} \\
\midrule
$\mathcal{C}_1$ & 3 & 5 & 40 & 506 & 0 & 3 & 557 & 0.908 & 0.505 \\
$\mathcal{C}_2$ & 17 & 7 & 280 & 29 & 0 & 0 & 333 & 0.841 & 0.778 \\
$\mathcal{C}_3$ & 2 & 1 & 1 & 7 & 0 & 2 & 13 & 0.538 & 1.761 \\
$\mathcal{C}_4$ & 7 & 1 & 1 & 0 & 422 & 3 & 434 & 0.972 & 0.239 \\
$\mathcal{C}_5$ & 40 & 180 & 20 & 10 & 2 & 3 & 255 & 0.706 & 1.341 \\
$\mathcal{C}_6$ & 4 & 505 & 6 & 0 & 0 & 4 & 519 & 0.973 & 0.208 \\
$\mathcal{C}_7$ & 423 & 3 & 12 & 3 & 0 & 1 & 442 & 0.957 & 0.320 \\
$\mathcal{C}_8$ & 105 & 3 & 5 & 1 & 0 & 0 & 114 & 0.921 & 0.449 \\
$\mathcal{C}_9$ & 60 & 17 & 85 & 8 & 0 & 367 & 537 & 0.683 & 1.258 \\
\midrule
\textbf{Globale} & & & & & & & \textbf{3204} & \textbf{0.887} & \textbf{0.672} \\
\bottomrule
\end{tabular}%
}
\caption{Matrice di contingenza e calcolo di entropia e purezza per il clustering a $K=9$ del dataset \textit{LA Documents} (Tan et al.~\cite{tan2018}).}
\label{tab:dm-la-documents}
\end{table}

\begin{noteblock}{Interpretazione Qualitativa dei Risultati}
Dall'ispezione analitica della Tabella~\ref{tab:dm-la-documents} emergono osservazioni cruciali:
\begin{enumerate}
    \item \textbf{Cluster ad altissima purezza}: il cluster $\mathcal{C}_4$ (422 articoli di Sport su 434, purezza $97.2\%$) e il cluster $\mathcal{C}_6$ (505 articoli di Finanza su 519, purezza $97.3\%$) manifestano una separazione semantica quasi perfetta. Le loro entropie sono inferiori a $0.25$;
    \item \textbf{Sotto-partizione naturale della classe Intrattenimento}: la classe \textit{Entertainment} è stata suddivisa dall'algoritmo tra il cluster $\mathcal{C}_7$ (423 articoli) e il cluster $\mathcal{C}_8$ (105 articoli). Sebbene tale classe sia frammentata su due cluster, entrambi i cluster risultano altamente puri ($95.7\%$ e $92.1\%$), a dimostrazione che l'algoritmo ha scoperto due sottotemi distinti all'interno della stessa macro-categoria;
    \item \textbf{Cluster misti complessi}: il cluster $\mathcal{C}_9$ racchiude 367 articoli di Tecnologia, ma è contaminato da 85 articoli Esteri e 60 di Intrattenimento, determinando un'entropia elevata ($1.258$).
\end{enumerate}
\end{noteblock}

% ====================================================================
% SEZIONE 10.9: MODELLI NULLI E INFERENZA STATISTICA
% ====================================================================
\section{Quadro Statistico Inferenziale e Modelli Nulli}
\label{sec:dm-statistical-null-models}

Come stabilito nella Sezione~\ref{sec:dm-validity-foundations}, un valore di validità (ad esempio un $\mathrm{SSE} = 24.5$ o una correlazione matriciale $\rho = 0.42$) è privo di significato epistemologico se non confrontato con il valore atteso sotto una distribuzione di probabilità di riferimento. L'approccio rigoroso della cluster validity si fonda sull'\textbf{inferenza statistica basata su modelli nulli}~\cite{tan2018, jain1988}.

\subsection{Formulazione dell'Ipotesi Nulla di Randomness}
Per stabilire la significatività statistica di un clustering, si formula formalmente l'ipotesi nulla:
\begin{equation}
\label{eq:dm-null-hypothesis}
H_0: \text{I dati sono privi di cluster e generati casualmente.}
\end{equation}

Nel contesto dei dati continui multidimensionali, l'ipotesi nulla standard modella le osservazioni come un processo di Poisson omogeneo generato uniformemente all'interno dell'iper-rettangolo o dell'iper-sfera che racchiude i dati osservati:
\begin{equation}
\mathbf{x} \sim \mathcal{U}\left(\prod_{j=1}^d [a_j, b_j]\right)
\end{equation}
oppure mediante una permutazione casuale indipendente degli attributi lungo ciascuna coordinata (preservando le distribuzioni marginali unidimensionali ma distruggendo la covarianza multivariata).

\subsection{Simulazione Monte Carlo e Distribuzione Empirica Sotto $H_0$}
Poiché la distribuzione analitica di metriche complesse come l'$\mathrm{SSE}$ o la correlazione matriciale non è derivabile in forma chiusa sotto $H_0$, si ricorre al \textbf{campionamento Monte Carlo}:
\begin{enumerate}
    \item Si generano $B$ dataset casuali indipendenti $\mathcal{D}^{(1)}, \mathcal{D}^{(2)}, \dots, \mathcal{D}^{(B)}$ (tipicamente $B = 500$ o $1\,000$) con la medesima cardinalità $N$ e dimensionalità $d$ del dataset reale sotto l'ipotesi nulla $H_0$;
    \item Si esegue l'algoritmo di clustering su ciascun dataset sintetico $\mathcal{D}^{(b)}$ con le medesime condizioni e iperparametri;
    \item Si calcola il valore della metrica statistica target $\theta^{(b)}$ (ad esempio l'$\mathrm{SSE}$ o la correlazione $\rho$) su ciascun dataset sintetico;
    \item Si costruisce la distribuzione empirica nulla $f(\theta \mid H_0)$.
\end{enumerate}

\subsection{Calcolo del $p$-Value Empirico}
Sia $\theta_{\mathrm{obs}}$ il valore della statistica calcolato sul dataset reale.
\begin{itemize}
    \item Per metriche in cui la bontà è associata alla minimizzazione (come l'$\mathrm{SSE}$), il $p$-value empirico è la proporzione di dataset casuali che hanno prodotto un errore uguale o inferiore a quello osservato:
    \begin{equation}
    \label{eq:dm-p-value-minimization}
    p\text{-value} = \frac{1}{B} \sum_{b=1}^B \mathbb{I}\left(\theta^{(b)} \le \theta_{\mathrm{obs}}\right)
    \end{equation}
    \item Per metriche in cui la qualità è associata alla massimizzazione (come la correlazione matriciale $\rho$ o la Silhouette):
    \begin{equation}
    \label{eq:dm-p-value-maximization}
    p\text{-value} = \frac{1}{B} \sum_{b=1}^B \mathbb{I}\left(\theta^{(b)} \ge \theta_{\mathrm{obs}}\right)
    \end{equation}
\end{itemize}

\begin{figure}[htbp]
\centering
\resizebox{\linewidth}{!}{%
\begin{tikzpicture}[
    scale=0.92,
    every node/.style={font=\footnotesize}
]
    % Decision zone horizontal braces below the axis (y = -0.75)
    % Rejection Region brace
    \draw[decorate, decoration={brace, mirror, amplitude=5pt}, crimson!80!black, thick] 
        (0.5, -0.75) -- (6.8, -0.75) 
        node[midway, below=5pt, font=\scriptsize, align=center, text=crimson!90!black] 
        {\textbf{Regione di Rigetto di $H_0$} ($\theta \le \theta_{\mathrm{crit}}$)\\[1.5pt]Evidenza di cluster reali ($p < \alpha$)};

    % Non-rejection Region brace
    \draw[decorate, decoration={brace, mirror, amplitude=5pt}, navy!70!black, thick] 
        (6.8, -0.75) -- (12.4, -0.75) 
        node[midway, below=5pt, font=\scriptsize, align=center, text=navy!85!black] 
        {\textbf{Regione di Non-Rigetto} ($\theta > \theta_{\mathrm{crit}}$)\\[1.5pt]Compatibile con rumore ($p \ge \alpha$)};

    % Shaded areas under the null distribution curve
    % Critical tail (rejection tail under H0)
    \fill[crimson!25, domain=6.0:6.8, samples=30] 
        (6.0, 0) -- plot (\x, {3.2 * exp(-0.5 * ((\x - 9.2)/1.0)^2)}) -- (6.8, 0) -- cycle;

    % Acceptance area under H0
    \fill[navy!12, domain=6.8:12.2, samples=60] 
        (6.8, 0) -- plot (\x, {3.2 * exp(-0.5 * ((\x - 9.2)/1.0)^2)}) -- (12.2, 0) -- cycle;

    % Outline of the null distribution bell curve
    \draw[navy!85!black, thick, domain=6.0:12.2, samples=80] 
        plot (\x, {3.2 * exp(-0.5 * ((\x - 9.2)/1.0)^2)});

    % Axes
    \draw[-Stealth, thick] (0, 0) -- (12.8, 0) node[above left=2pt, font=\footnotesize] {Valore Statistico $\theta$ ($\mathrm{SSE}$)};
    \draw[-Stealth, thick] (0, 0) -- (0, 4.6) node[above=2pt, font=\footnotesize] {Frequenza Relativa};

    % Null distribution labels placed clearly above the bell curve peak
    \node[navy!90!black, font=\bfseries, align=center] at (9.5, 3.9) {Distribuzione Nulla $f(\theta \mid H_0)$\\[2pt]\scriptsize (500 simulazioni Monte Carlo)};

    % Mean of null distribution tick
    \draw[gray!60, thin] (9.2, 0.1) -- (9.2, -0.1);
    \node[below=2pt, gray!70!black, font=\scriptsize] at (9.2, 0) {$\mu_{H_0}$};

    % Critical boundary threshold (x = 6.8)
    \draw[dashed, darkgreen!80!black, thick] (6.8, 0) -- (6.8, 1.75);
    \fill[darkgreen!80!black] (6.8, 0) circle (2pt);
    \node[below=2pt, darkgreen!80!black, font=\scriptsize\bfseries] at (6.8, 0) {$\theta_{\mathrm{crit}}$};
    \node[draw=darkgreen!80!black, fill=darkgreen!8, rounded corners=3pt, font=\scriptsize\bfseries, align=center, text=darkgreen!90!black, inner sep=3.5pt] 
        at (6.8, 2.25) {Soglia Critica $\alpha = 0.01$\\[1pt]\normalfont\scriptsize ($z_{\mathrm{crit}} \approx -2.33$)};

    % Critical tail callout (x = 4.6)
    \draw[-Stealth, darkgreen!80!black, thin] (4.9, 1.1) -- (6.3, 0.2);
    \node[darkgreen!90!black, font=\scriptsize, align=center] at (4.5, 1.4) {Coda critica\\($\alpha = 0.01$)};

    % Observed value (x = 2.0)
    \fill[crimson] (2.0, 0) circle (2.5pt);
    \node[below=2pt, crimson!90!black, font=\scriptsize\bfseries] at (2.0, 0) {$\theta_{\mathrm{obs}}$};
    \draw[crimson, very thick, -Stealth] (2.0, 2.8) -- (2.0, 0.2);
    \node[draw=crimson, fill=crimson!8, rounded corners=3pt, font=\footnotesize, align=center, text=crimson!90!black, inner sep=4pt] 
        at (2.0, 3.4) {\textbf{Valore Osservato} $\theta_{\mathrm{obs}}$\\[2pt]\scriptsize Dati reali ($p\text{-value} < 0.001$)};

\end{tikzpicture}%
}
\caption{Quadro decisionale dell'ipotesi nulla: se l'errore $\mathrm{SSE}$ reale cade marcatamente a sinistra della distribuzione generata casualmente, l'ipotesi di casualità viene rigettata con significatività statistica rigorosa.}
\label{fig:dm-null-model-distribution}
\end{figure}

Come sintetizzato graficamente nella Figura~\ref{fig:dm-null-model-distribution}, se il valore empirico osservato $\theta_{\mathrm{obs}}$ risiede ben al di fuori della coda di probabilità della distribuzione nulla ($p\text{-value} < \alpha$, tipicamente $\alpha = 0.01$), l'ipotesi di uniformità viene rigettata: la partizione non è un mero artefatto algoritmico ma riflette una discontinuità reale nei dati.

% ====================================================================
% SEZIONE 10.10: LABORATORIO PYTHON APPLICATIVO
% ====================================================================
\section{Laboratorio Python: Metriche di Validazione e Selezione di \texorpdfstring{$K$}{K}}
\label{sec:dm-python-lab}

In questa sezione finale vengono presentate le metodologie operative implementate nell'attività laboratoriale di data mining (desunte dal notebook ufficiale del corso \texttt{3\_clustering.ipynb}~\cite{tan2018}), evidenziando l'impiego del package \texttt{scikit-learn} per la quantificazione delle metriche interne ed esterne e la stima del numero ottimale di cluster.

\begin{noteblock}{Ambiente Interattivo e Riproducibilità del Laboratorio}
    Tutte le routine computazionali, le metriche di validazione e le procedure di model selection presentate in questa sezione sono direttamente eseguibili e riproducibili tramite l'ambiente Jupyter ufficiale del corso:
    \begin{itemize}[noitemsep]
        \item Notebook interattivo: \href{http://localhost:8888/lab/tree/notebooks/3_clustering.ipynb}{\textbf{Notebook 3 (Clustering)}} (file \texttt{Data-Mining/notebooks/3\_clustering.ipynb});
        \item Avvio rapido con Docker: \texttt{docker compose up jupyter} (oppure tramite script Windows \texttt{scripts/start-jupyter.cmd});
        \item Interfaccia web: \href{http://localhost:8888/lab}{\texttt{http://localhost:8888/lab}} (accesso diretto senza token o password).
    \end{itemize}
\end{noteblock}

\begin{remark}[Riferimento al Notebook Interattivo]
    Il codice delle sottosezioni seguenti mappa direttamente le celle della sezione \textit{Metrics to Evaluate Clustering} e \textit{K parameter estimation} del notebook interattivo \href{http://localhost:8888/lab/tree/notebooks/3_clustering.ipynb}{\textbf{3\_clustering.ipynb}}:
    \begin{itemize}[noitemsep]
        \item \textbf{Metriche interne} ($\mathrm{SSE}$, Davies-Bouldin, Silhouette): celle 33--34;
        \item \textbf{Metodo del gomito} (scansione di $K \in [2, 40]$ ed evidenziazione di $K^*$): celle 38--40;
        \item \textbf{Validazione esterna} (matrice di contingenza e profilazione con \texttt{day}): celle 35--37.
    \end{itemize}
\end{remark}

\subsection{Implementazione Metriche Interne in Scikit-Learn}
Nel codice di laboratorio (cella 34 del \href{http://localhost:8888/lab/tree/notebooks/3_clustering.ipynb}{\textbf{Notebook 3}}), dopo aver addestrato un modello $K$-Means sulla matrice pre-processata dei dati numerici standardizzati $X$, le metriche di compattezza ($\mathrm{SSE}$), separazione (Davies-Bouldin) e concordanza (Silhouette) vengono estratte in modo conciso:

\begin{lstlisting}[style=code, language=Python, caption={Calcolo congiunto di $\mathrm{SSE}$, Davies-Bouldin e Silhouette Score in Python.}]
from sklearn.cluster import KMeans
from sklearn import metrics
from sklearn.metrics import silhouette_score

# Addestramento di K-Means con K fissato
kmeans = KMeans(n_clusters=4, n_init=10, max_iter=300, random_state=42)
kmeans.fit(X)

# 1. Coesione intra-cluster: Inerzia (SSE / WSS)
# Attributo built-in che calcola la somma delle distanze quadratiche dai centroidi
sse = kmeans.inertia_
print(f'SSE (Inertia / WSS): {sse:.3f}')

# 2. Separazione inter-cluster: Davies-Bouldin Score
# Valuta il rapporto di similarita peggiore; ottimale vicino a 0
db_score = metrics.davies_bouldin_score(X, kmeans.labels_)
print(f'Davies-Bouldin Score: {db_score:.3f}')

# 3. Silhouette Score globale: media armonica tra coesione e separazione
# Spettro [-1, 1]; ottimale vicino a +1
sil_score = silhouette_score(X, kmeans.labels_)
print(f'Silhouette Score: {sil_score:.3f}')
\end{lstlisting}

\subsection{Stima del Parametro \texorpdfstring{$K$}{K} Ottimale: Il Metodo del Gomito (Elbow/Knee)}
Come formalizzato nel Capitolo~\ref{chap:dm-kmeans-hierarchical} (\autoref{sec:dm-kmeans-elbow-method}), l'identificazione euristica di $K^*$ si basa sull'individuazione del punto di flesso (\textit{knee point}) in cui il guadagno marginale di coesione interna subisce una marcata decelerazione.

Il listato~\ref{lst:dm-elbow-method} (desunto dalle celle 38--40 del \href{http://localhost:8888/lab/tree/notebooks/3_clustering.ipynb}{\textbf{Notebook 3}}) illustra la scansione programmatica dell'intervallo $K \in [2, 40]$ e il tracciamento della curva per individuare visivamente la transizione a gomito:

\begin{lstlisting}[style=code, language=Python, label={lst:dm-elbow-method}, caption={Scansione sistematica e tracciamento della curva del gomito per la selezione di $K^*$.}]
import matplotlib.pyplot as plt

sse_list = []
max_k = 40

# Iterazione sull'intervallo di complessita
for k in range(2, max_k + 1):
    km = KMeans(n_clusters=k, n_init=10, max_iter=100, random_state=42)
    km.fit(X)
    sse_list.append(km.inertia_)

# Identificazione del punto critico di ginocchio
k_star = 10
sse_star = sse_list[k_star - 2]  # Allineamento indice base-zero

# Visualizzazione grafica della curva del gomito
plt.figure(figsize=(9, 5))
plt.plot(range(2, max_k + 1), sse_list, marker='o', color='navy', lw=2)
plt.ylabel('SSE (Inertia)', fontsize=14)
plt.xlabel('Numero di Cluster (K)', fontsize=14)
plt.title('Metodo del Gomito per la Stima di K*', fontsize=16)

# Linee tratteggiate di ancoraggio sul punto ottimale
plt.vlines(x=k_star, ymin=0, ymax=sse_star, colors='crimson', linestyles='dashed', lw=1.5)
plt.hlines(y=sse_star, xmin=0, xmax=k_star, colors='crimson', linestyles='dashed', lw=1.5)
plt.scatter(k_star, sse_star, color='crimson', s=90, zorder=5, label=f'Ginocchio K*={k_star}')

plt.legend(fontsize=12)
plt.grid(True, linestyle=':', alpha=0.6)
plt.show()
\end{lstlisting}

\subsection{Validazione Esterna tramite Tabulazioni Incrociate (Crosstab)}
Nel laboratorio (celle 35--37 del \href{http://localhost:8888/lab/tree/notebooks/3_clustering.ipynb}{\textbf{Notebook 3}}), per esaminare se i cluster numerici scoperti corrispondano a comportamenti differenziati rispetto a variabili categoriche esterne non utilizzate nel clustering (come il giorno della settimana \texttt{day} o la fascia oraria \texttt{time} nel dataset delle mance \texttt{tips}), si impiega la funzione \texttt{pd.crosstab}:

\begin{lstlisting}[style=code, language=Python, caption={Analisi di contingenza empirica e profilazione esterna dei cluster in Pandas.}]
import pandas as pd

# Creazione della matrice di contingenza tra etichette di cluster e variabile esterna
cross_df = pd.crosstab(kmeans.labels_, df_cat['day'])
print("Matrice di Contingenza Cluster vs Giorno:")
print(cross_df)

# Visualizzazione della distribuzione per cluster
cross_df.plot(kind='bar', stacked=False, figsize=(10, 5), colormap='viridis')
plt.title('Distribuzione della variabile esterna "day" per cluster')
plt.xlabel('Cluster ID')
plt.ylabel('Conteggio Istanza')
plt.grid(axis='y', linestyle=':', alpha=0.7)
plt.show()
\end{lstlisting}

La matrice di contingenza consente all'analista di calcolare immediatamente la purezza e l'entropia condizionata del partizionamento rispetto alle etichette categoriche, completando in tal modo il ciclo di valutazione quantitativa della validità.
